%% Exemplo de uso do sistema de folha de aprovação do PPGT
%% Este arquivo demonstra como usar os novos comandos para definir membros da banca

\documentclass[mestrado]{UnB-PPGT}

%% ===== METADADOS BÁSICOS =====
\titulo{Análise de Sistemas de Transporte Urbano Sustentável}
\subtitulo{Uma Abordagem Integrada para Mobilidade Urbana}
\autor{João}{Silva}

%% ===== ORIENTAÇÃO =====
\orientador{Prof. Dr. Maria Santos}{Universidade de Brasília}
\coorientador{Prof. Dr. Carlos Oliveira}{Universidade Federal do Rio de Janeiro}

%% ===== DATA DE DEFESA =====
\datadefesa{15}{12}{2024}

%% ===== COORDENADOR DO PROGRAMA =====
\coordenador{Prof. Dr. Ana Costa}{Universidade de Brasília}

%% ===== MEMBROS DA BANCA =====
%% Para dissertação de mestrado: mínimo 3 membros (orientador + 2 examinadores)
%% Para tese de doutorado: mínimo 5 membros (orientador + 4 examinadores)

% Comando básico para membro da banca (assume "Examinador")
\membrobanca{Prof. Dr. Pedro Almeida}{Universidade de São Paulo}
\membrobanca{Prof. Dr. Lucia Ferreira}{Universidade Federal de Minas Gerais}

% Comandos específicos para diferentes tipos de membros
\examinadorinterno{Prof. Dr. Roberto Lima}{Universidade de Brasília}
\examinadorexterno{Prof. Dr. Ana Souza}{Universidade Federal do Rio de Janeiro}

% Comando para membro com função específica
\membroespecial{Prof. Dr. Carlos Mendes}{Instituto de Pesquisa}{Especialista Convidado}

%% ===== PALAVRAS-CHAVE =====
\palavraschave{transporte urbano, mobilidade sustentável, sistemas integrados}
\keywords{urban transport, sustainable mobility, integrated systems}

%% ===== NÚMERO DA PUBLICAÇÃO (OPCIONAL) =====
\numeropublicacao{PPGT.TD-001/2024}

\begin{document}

%% ===== PÁGINAS PRÉ-TEXTUAIS =====
\pretextual

% A folha de aprovação será gerada automaticamente pelo \maketitle
% ou pode ser chamada diretamente com \folhaaprovacao

%% ===== CONTEÚDO TEXTUAL =====
\textual

\chapter{Introdução}

Este documento demonstra o uso do sistema aprimorado de folha de aprovação 
do modelo LaTeX para o PPGT-UnB.

\section{Novos Comandos Implementados}

\subsection{Comandos Básicos para Membros da Banca}

\begin{itemize}
\item \texttt{\textbackslash membrobanca\{Nome\}\{Instituição\}} - Define um membro da banca como "Examinador"
\item \texttt{\textbackslash examinadorinterno\{Nome\}\{Instituição\}} - Define um examinador interno
\item \texttt{\textbackslash examinadorexterno\{Nome\}\{Instituição\}} - Define um examinador externo
\item \texttt{\textbackslash membroespecial\{Nome\}\{Instituição\}\{Função\}} - Define um membro com função específica
\end{itemize}

\subsection{Características do Sistema}

\begin{itemize}
\item Validação automática do número mínimo de membros
\item Geração automática de espaços para assinaturas
\item Layout profissional seguindo padrões do PPGT
\item Suporte para diferentes números de membros
\item Formatação específica para mestrado e doutorado
\end{itemize}

\chapter{Conclusão}

O sistema implementado atende aos requisitos 3.3 e 2.4 especificados:
\begin{itemize}
\item Gera automaticamente a folha de aprovação com espaços para assinaturas da banca
\item Cria automaticamente a folha de aprovação quando membros da banca são definidos
\end{itemize}

\end{document}